\documentclass[../Main.tex]{subfiles}
\begin{document}
  \chapter{2019-2020学年微积分（一）（上）期中考试}

  \section{基本计算题(每小题 6 分，共 60 分)}
  \begin{enumerate}
    \item 求极限$\lim_{n\to \infty}\sqrt[n]{2\sin^{2} n+3\cos^{2} n}$.
      \vspace{10em}

    \item 当$x\to1$时，求$\sqrt[3]{x}-1$关于$x-1$的主部.
      \vspace{12em}

    \item 计算极限$l=\lim_{x\to0}\frac{\sin x-x\cos x}{x\ln\cos x}$.
      \vspace{12em}

    \item 计算极限$l=\lim_{x\to0}\left(\mathrm{e}^{x}-\sin x\right)^{\frac{1}{x^{2}}}$.
      \vspace{11em}

    \item 设$f(x)=x^{\tan x}$.求$f'\left(\frac{\pi}{4}\right)$.
      \vspace{9em}

    \item 求曲线$\cos x+\ln(y-x)=y^{2}$在点$(0,1)$处的切线方程.
      \vspace{10em}

    \item 设函数$f(x)$二阶可导，求$y=f(\sin x)$的二阶导数.
      \vspace{10em}

    \item 设$y=x\ln x$，求$y^{(6)}$.
      \vspace{10em}

    \item 求$a,b$的值使得函数$f(x)=
      \begin{cases}
        \mathrm{e}^{ax}, & x<0,   \\
        b-\sin 2x,       & x\geq0
      \end{cases}$在$x=0$处连续并且可导.
      \vspace{8em}

    \item 设函数$y=y(x)$由方程$\begin{cases}
        x=t^{2}+2t, \\
        y=t+\ln t
      \end{cases}$确定，求$\frac{\mathrm{d}y}{\mathrm{d}x}$和$\frac{\mathrm{d}^{2}
      y}{\mathrm{d}x^{2}}$.
      \vspace{8em}
  \end{enumerate}

  \section{综合题(每小题 6 分，共 30 分)}
  \begin{enumerate}
    \item[11.] 设函数$f(x)=\mathrm{e}^{\frac{1}{x^{2}}}\arctan\frac{x-1}{x+1}$.指出$f
      (x)$的间断点，并判断其类型.
      \vspace{8em}

    \item[12.] 设$f(x)=x^{n}+a_{1} x^{n-1}+\cdots+a_{n-1}x+a_{n}$为$n$次实系数多项式.证明当$n$为奇数时，方程$f
      (x)=0$至少有一个实根.
      \vspace{8em}

    \item[13.] 设$f(x)$在$x=2$处可导，$f(2)\neq0$，求$\lim_{x\to\infty}\left[\frac{f\left(2+\frac{1}{x}\right)}{f(2)}
      \right]^{x}$.
      \vspace{9em}

    \item[14.] 设一个雪球以$2\mathrm{cm}^{3}/\mathrm{min}$的速度融化.设雪球在融化过程中始终保持球形.求当雪球半径为$1
      0\mathrm{cm}$的时候半径变化的速率.
      \vspace{10em}

    \item[15.] 写出函数$f(x)=x\cos x$带Peano余项的五阶Maclaurin公式.
      \vspace{12em}
  \end{enumerate}

  \section{证明题(每小题 5 分，共 10 分)}
  \begin{enumerate}
    \item[16.] 设数列$\{x_{n}\}$由递推公式$x_{1}=\sqrt{6}$，$x_{n+1}=\sqrt{6+x_{n}}$给出.证明$\lim
      _{n\to\infty}x_{n}$存在并求极限值.
      \vspace{12em}

    \item[17.] 设$0<a<b$，函数$f(x)$在闭区间$[a,b]$上连续，在开区间$(a,b)$内可导.证明存在$\xi
      \in(a,b)$使得
      \[
        \xi f'(\xi)-f(\xi)=\frac{af(b)-bf(a)}{b-a}.
      \]
  \end{enumerate}

  \chapter{2019-2020学年微积分（一）（上）期中考试参考答案}
  \section{基本计算题(每小题 6 分，共 60 分)}
  \begin{enumerate}
    \item \textbf{Solution}. 显然$\sqrt[n]{2}\leq\sqrt[n]{2\sin^{2} n+3\cos^{2}
      n}\leq\sqrt[n]{3}$.

      由夹逼定理知$\lim_{n\to\infty}\sqrt[n]{2\sin^{2} n+3\cos^{2} n}=1$.

    \item \textbf{Solution}.
      \[
        \sqrt[3]{x}-1=(1+x-1)^{\frac{1}{3}}-1\sim\frac{1}{3}(x-1).
      \]

      故主部为$\frac{1}{3}(x-1)$.

    \item \textbf{Solution}.
      \[
        \begin{aligned}
          l= & \lim_{x\to0}\frac{\sin x-x\cos x}{-\frac{1}{2}x^3} \\
          =  & -2\lim_{x\to0}\frac{x\sin x}{3x^2}=-\frac{2}{3}.
        \end{aligned}
      \]

    \item \textbf{Solution}.
      \[
        \begin{aligned}
          l= & \lim_{x\to0}\left(\mathrm{e}^{x}-\sin x\right)^{\frac{1}{x^2}}=\mathrm{e}^{\lim\limits_{x\to0} \frac{1}{x^2}\ln\left(\mathrm{e}^x-\sin x\right)}                 \\
          =  & \mathrm{e}^{\lim\limits_{x\to0}\frac{\mathrm{e}^x-\cos x}{\mathrm{e}^x-\sin x}\cdot\frac{1}{2x}}= \mathrm{e}^{\lim\limits_{x\to0}\frac{\mathrm{e}^x-\cos x}{2x}} \\
          =  & \mathrm{e}^{\lim\limits_{x\to0}\frac{\mathrm{e}^x+\sin x}{2}}=\sqrt{\mathrm{e}}.
        \end{aligned}
      \]

    \item \textbf{Solution}. $\ln f(x)=\tan x\ln x$，

      所以$f'(x)=f(x)\cdot\left(\tan x\ln x\right)'=x^{\tan x}\left(\sec^{2} x \ln
      x + \frac{\tan x}{x}\right)$，

      代入$x=\frac{\pi}{4}$得$f'\left(\frac{\pi}{4}\right)=1+\frac{\pi}{2}\ln\frac{\pi}{4}$.

    \item \textbf{Solution}. 方程两边对$x$求导得到
      \[
        -\sin x+\frac{1}{y-x}(y'-1)=2yy',
      \]

      代入$x=0,y=1$解得$y'(0)=-1$，所以切线方程为$x+y=1$.

    \item \textbf{Solution}. $y'=f'(\sin x)\cos x$，

      $y''=f''(\sin x)\cos^{2} x-f'(\sin x)\sin x$.

    \item \textbf{Solution}.

      \textbf{法一}. 根据Leibniz法则得到
      \[
        \begin{aligned}
          y^{(6)}= & x\left(\ln x\right)^{(6)}+6\left(\ln x\right)^{(5)}             \\
          =        & x\left(\frac{1}{x}\right)^{(5)}+6\left(\frac{1}{x}\right)^{(4)} \\
          =        & \frac{24}{x^5}.
        \end{aligned}
      \]

      \textbf{法二}.
      \[
        y^{(6)}=\left(1+\ln x\right)^{(5)}=\left(\frac{1}{x}\right)^{(4)}=\frac{24}{x^{5}}
        .
      \]

    \item \textbf{Solution}. 直接得到$f\left(0^{+}\right)=b$，$f\left(0^{-}\right
      )=1$，所以$b=1$.

      求导得到$f'_{+}(0)=-2$，$f'_{-}(0)=a$，所以$a=-2$.

    \item \textbf{Solution}. $\frac{\mathrm{d}y}{\mathrm{d}x}=\frac{\frac{\mathrm{d}y}{\mathrm{d}t}}{\frac{\mathrm{d}x}{\mathrm{d}t}}
      =\frac{1+\frac{1}{t}}{2t+2}=\frac{1}{2t}$.

      $\frac{\mathrm{d}^{2} y}{\mathrm{d}x^{2}}=\frac{\mathrm{d}}{\mathrm{d}t}\left
      (\frac{\mathrm{d}y}{\mathrm{d}x}\right)\cdot\frac{\mathrm{d}t}{\mathrm{d}x}
      =-\frac{1}{2t^{2}}\cdot\frac{1}{2t+2}=-\frac{1}{4t^{2}(t+1)}$.
  \end{enumerate}

  \section{综合题(每小题 6 分，共 30 分)}
  \begin{enumerate}
    \item[11.] \textbf{Solution}. 函数$f(x)$的间断点为$0,-1$. 因
      \[
        \lim_{x\to0}f(x)=-\infty,\quad f(-1^{+})=-\frac{\pi\mathrm{e}}{2},\quad f
        (-1^{-})=\frac{\pi\mathrm{e}}{2},
      \]

      所以$x=0$是无穷间断点，$x=-1$是跳跃间断点.

    \item[12.] \textbf{Solution}. 当$n$为奇数时，显然有
      \[
        f\left(+\infty\right)=+\infty,\quad f\left(-\infty\right)=-\infty.
      \]

      于是存在$x_{1},x_{2}$使得$f(x_{1})<0<f(x_{2})$.

      函数$f(x)$显然连续，由介值定理知，存在$x_{0}\in(x_{1},x_{2})$使得$f(x_{0})=
      0$，即方程$f(x)=0$至少有一个实根.

    \item[13.] \textbf{Solution}.
      \[
        \begin{aligned}
          \lim_{x\to\infty}\left[\frac{f\left(2+\frac{1}{x}\right)}{f(2)}\right]^{x} = & \mathrm{e}^{\lim\limits_{x\to\infty}x\ln\frac{f\left(2+\frac{1}{x}\right)}{f(2)}}                       \\
          =                                                                            & \mathrm{e}^{\lim\limits_{x\to\infty}x\ln\left(1+\frac{f\left(2+\frac{1}{x}\right)-f(2)}{f(2)}\right)}   \\
          =                                                                            & \mathrm{e}^{\frac{1}{f(2)}\lim\limits_{x\to\infty}x\left[f\left(2+\frac{1}{x}\right)-f(2)\right]}       \\
          =                                                                            & \mathrm{e}^{\frac{1}{f(2)}\lim\limits_{x\to\infty}\frac{f\left(2+\frac{1}{x}\right)-f(2)}{\frac{1}{x}}} \\
          =                                                                            & \mathrm{e}^{\frac{f'(2)}{f(2)}}.
        \end{aligned}
      \]

    \item[14.] \textbf{Solution}. 根据几何关系得到
      \[
        V=\frac{4\pi r^{3}}{3}.
      \]

      两边对$t$求导得到
      \[
        \frac{\mathrm{d}V}{\mathrm{d}t}=4\pi r^{2}\cdot\frac{\mathrm{d}r}{\mathrm{d}t}
        .
      \]

      代入$\frac{\mathrm{d}V}{\mathrm{d}t}=-2$以及$r=10$解得$\frac{\mathrm{d}r}{\mathrm{d}t}
      =-\frac{1}{200\pi}$.

      即半径以$\frac{1}{200\pi}\mathrm{cm}/\mathrm{min}$的速度减小.

    \item[15.] \textbf{Solution}. 先写出
      \[
        \cos x=1-\frac{1}{2}x^{2}+\frac{1}{24}x^{4}+o(x^{4}),
      \]

      进而
      \[
        f(x)=x\cos x=x-\frac{1}{2}x^{3}+\frac{1}{24}x^{5}+o(x^{5}).
      \]
  \end{enumerate}

  \section{证明题(每小题 5 分，共 10 分)}
  \begin{enumerate}
    \item[16.] \textbf{Proof}.

      显然$x_{2}>x_{1}$，设$x_{k+1}>x_{k}$，由$x_{k+2}^{2}-x_{k+1}^{2}=6+x_{k+1}-
      6-x_{k}=x_{k+1}-x_{k}>0$知$x_{k+2}>x_{k+1}$.

      由数学归纳法知$\{x_{n}\}$严格单调增加.

      另外显然$x_{1}<3$，设$x_{k}<3$，由$x_{k+1}=\sqrt{6+x_{k}}<\sqrt{9}=3$知$x_{k+1}
      <3$.

      由数学归纳法知$\{x_{n}\}$有上界3. 根据单调有界收敛准则知数列$\{x_{n}\}$极限存在.

      设$\lim_{n\to\infty}x_{n}=L$，对$x_{n+1}=\sqrt{6+x_{n}}$两边取极限得到$L=\sqrt{6+L}$，解得$L
      =3$.

    \item[17.] \textbf{Proof}. 显然函数$F(x)=\frac{f(x)}{x}$和$G(x)=\frac{1}{x}$在闭区间$[
      a,b]$上连续，在开区间$(a,b)$内可导，并且$G(x)$的导数不为0.

      由Cauchy中值定理知存在$\xi\in(a,b)$使得
      \[
        \frac{af(b)-bf(a)}{a-b}=\frac{F(b)-F(a)}{G(b)-G(a)}=\frac{F'(\xi)}{G'(\xi)}
        =\xi f'(\xi)-f(\xi).
      \]

      两边变号即得所证.
  \end{enumerate}
\end{document}
